\documentclass[12pt]{article}
\usepackage[utf8]{inputenc}
\usepackage[spanish]{babel}
\usepackage{amsmath, amssymb, amsfonts}
\usepackage{geometry}
\geometry{margin=1in}

\title{Ejercicios: Modelo de Probabilidad Lineal}
%\author{}
\date{}

\begin{document}

\maketitle



{\bf Ejercicio 1:} Sea una variable dependiente dicotómica \( Y_i \) tal que:
\[
Y_i =
\begin{cases}
1 & \text{si ocurre el evento de interés,} \\
0 & \text{en caso contrario.}
\end{cases}
\]
Se asume que \( Y_i \) sigue una distribución Bernoulli con parámetro \( p_i \), es decir:
\[
Y_i | X_i \sim \text{Bernoulli}(p_i),
\]
donde \( p_i = P(Y_i = 1|X_i) \).

En el \textbf{Modelo de Probabilidad Lineal (MPL)}, la probabilidad condicional se especifica de forma lineal:
\[
E(Y_i|X_i) = p_i = X_i'\beta,
\]
donde \( X_i'\beta = \beta_0 + \beta_1 X_{1i} + \cdots + \beta_k X_{ki} \).

---


El error se define como:
\[
u_i = Y_i - E(Y_i|X_i) = Y_i - p_i.
\]

La varianza condicional del error es:
\[
\operatorname{Var}(u_i|X_i) = \operatorname{Var}(Y_i|X_i).
\]
Dado que \( Y_i|X_i \) sigue una Bernoulli(\(p_i\)), su varianza es:
\[
\operatorname{Var}(Y_i|X_i) = p_i(1 - p_i).
\]
Por tanto:
\[
\boxed{\operatorname{Var}(u_i|X_i) = p_i(1 - p_i)}.
\]



Dado que en el MPL la probabilidad se modela linealmente como:
\[
p_i = X_i'\beta,
\]
se tiene que:
\[
\operatorname{Var}(u_i|X_i) = X_i'\beta(1 - X_i'\beta).
\]
Claramente, esta varianza depende de \( X_i \).

La varianza condicional del error no es constante, ya que depende de los valores de las variables explicativas \( X_i \). Por lo tanto, el modelo de probabilidad lineal presenta:

\[
\boxed{\text{Heterocedasticidad: } \operatorname{Var}(u_i|X_i) \neq \sigma^2.}
\]


{\bf Ejercicio 2:} 
Suponga que usted cuenta con información sobre la \textit{felicidad} de las personas (\(B_i\)), la cual es una variable dummy que puede tomar sólo dos valores: \(0\) si el individuo se clasifica como “no feliz” y \(1\) si el individuo dice “ser feliz”. Por otro lado, asuma que tiene información, para estos mismos individuos, sobre un conjunto de variables que afectan la felicidad. En concreto, se conoce la edad (\(E_i\)), género (\(M_i\): mujer=1), años de escolaridad (\(S_i\)), ingreso familiar (\(I_i\), en miles de pesos al año) y estado de salud (\(H_i\): buena salud=1). Con base en esto se plantea el siguiente modelo econométrico:
\[
B_i \;=\; \alpha_1 + \alpha_2 E_i + \alpha_3 M_i + \alpha_4 S_i + \alpha_5 I_i + \alpha_6 H_i + e_i
\tag{1}
\]
donde \(e_i\) es un shock estocástico que captura factores no observables que afectan la felicidad.

\begin{enumerate}
\item[1.] ¿Es posible estimar el modelo (1) por Mínimos Cuadrados Ordinarios (MCO)? Justifique.
\item[2.] En el contexto de MCO, ¿cuál es la probabilidad de que un individuo sea feliz?
\item[3.] ¿Cuál es el efecto marginal del ingreso familiar sobre la probabilidad de ser feliz?
\item[4.] Mencione y explique dos críticas a la estimación por MCO del modelo (1). 
\end{enumerate}

\bigskip
\hrule
\bigskip



\paragraph{1)} 
Sí, es \textit{posible} correr MCO sobre (1); esto se conoce como \textbf{Modelo de Probabilidad Lineal} (MPL). Sin embargo, presenta problemas teóricos y de inferencia:

\begin{itemize}
\item \textbf{Heterocedasticidad inherente.} Como \(B_i\in\{0,1\}\) y \(p_i\equiv \mathbb{E}[B_i|X_i]=X_i'\alpha\), entonces
\[
u_i=B_i-p_i, \qquad \mathbb{V}(u_i|X_i)=\mathbb{V}(B_i|X_i)=p_i(1-p_i),
\]
la cual depende de \(X_i\) (\(p_i=X_i'\alpha\)), violando la homocedasticidad de MCO.
\item \textbf{Predicciones fuera de \([0,1]\).} La forma lineal puede entregar \(\hat p_i=\hat B_i\) menores que 0 o mayores que 1.
\item \textbf{Forma funcional inadecuada.} La relación entre \(X\) y una probabilidad suele ser no lineal; el efecto marginal constante de MCO es poco realista.
\item \textbf{Inferencia no válida sin corrección.} Los errores estándar clásicos de MCO son inconsistentes bajo heterocedasticidad; se requieren errores robustos (White).
\end{itemize}



\paragraph{2) Probabilidad de ser feliz bajo MCO (MPL).}
En el MPL la probabilidad condicional es la esperanza condicional:
\[
\mathbb{P}(B_i=1\,|\,X_i) \;=\; \mathbb{E}[B_i|X_i] \;=\; X_i'\alpha 
\;=\; \alpha_1 + \alpha_2 E_i + \alpha_3 M_i + \alpha_4 S_i + \alpha_5 I_i + \alpha_6 H_i.
\]
Empíricamente, \(\widehat{\mathbb{P}}(B_i=1|X_i)=X_i'\hat\alpha\) (eventualmente truncada a \([0,1]\) para interpretación).

\paragraph{3) Efecto marginal del ingreso en el MPL.}
En el modelo lineal, el efecto marginal es \textbf{constante} e igual al coeficiente:
\[
\frac{\partial \mathbb{P}(B_i=1|X_i)}{\partial I_i} \;=\; \alpha_5.
\]
(Si \(I_i\) está en miles de pesos/año; \(\alpha_5\) es el cambio en puntos de probabilidad por cada mil pesos adicionales).



\paragraph{4) Dos críticas a MCO y método sugerido.}
\begin{enumerate}
\item \textbf{Heterocedasticidad} \( \Rightarrow \) errores estándar mal calculados, pérdida de eficiencia y problemas con la inferencia. 
\item \textbf{Probabilidades fuera de rango y efectos marginales constantes,} lo que puede inducir interpretaciones erróneas.
\end{enumerate}

\end{document}





\end{document}